% =====================================================================
% VERSION ENSAE, REDUCTION SFE (phase 3, 1er octobre 2026) : resume.
% Copie a texte constant de 01_resume.tex.
% Ne changent que les titres et les renvois repares, listes dans
% reduction/phase3.md.
% Lu par main_ensae.tex seul : ne jamais le faire lire par main_institut.tex,
% qui garde chapitres/.
% =====================================================================

\chapter*{Résumé}
\addcontentsline{toc}{chapter}{Résumé}
% Un \chapter* ne met pas le titre courant a jour : sans cette ligne, la seconde page du
% resume portait << Remerciements >> en tete.
\markboth{Résumé}{}

% LONGUEUR CONTRAINTE PAR LE § 5.1.c DES RECOMMANDATIONS JURY : 150 A 250 MOTS, pour le
% resume francais comme pour l'abstract anglais, ces deux textes servant a l'indexation sur
% le site de l'Institut. Mesure du 23 septembre 2026 AVANT resserrage : 322 et 305 mots sur
% le PDF, soit +29 % et +22 %. Ce qui a ete retire est du DETAIL, jamais un resultat ni une
% reserve : le compte des sources qui echouent et le champ VERIS rempli une fois sur 10 591,
% et le mecanisme de la decroissance de la charge avec la taille. Tous trois vivent dans la
% note de synthese et dans l'executive summary, qui n'ont AUCUNE limite de longueur.
% NE PAS REGROSSIR CES DEUX BLOCS : recompter sur le PDF apres toute retouche.
% REECRITS LE 30 SEPTEMBRE 2026, 247 et 239 mots : la bande dite hors des deux etats, le
% chiffre d'entite en majorant, la seule dependance gouvernance-incidents, le groupe de tete
% des piliers, les deux canaux a bornes posees, et trois antitheses retirees. L'abstract
% passe en romain, comme le resume francais.
Le règlement européen DORA impose aux entités financières de maîtriser leurs risques informatiques.
Nous comparons le besoin de capital entre états de conformité supposés, puis examinons sa sensibilité
aux informations manquantes sur la contagion et les coûts.

Le modèle combine le nombre annuel d'incidents, leurs pertes et leur propagation entre les cinq
domaines de contrôle. Le besoin de capital est le quantile à $99{,}5\,\%$ de la perte annuelle,
avec addition des coûts par pilier touché. Les pénalités et les polices cyber vendues sont exclues.

Sur le périmètre financier étudié, le besoin passe de $6\,049$ à $20\,188$~M\euro{} entre les états
conforme et non conforme, soit un facteur $3{,}34$. Ces montants dépendent des paramètres de scénario
et de la règle de coût. Pour une entité de $20\,000$~M\$ d'actifs, la plage de $131{,}5$ à
$183{,}3$~M\euro{} reste une extrapolation d'ordre de grandeur.

Les sources ne permettent pas de calibrer la propagation. Nous conservons la force des liens
posée à dire d'expert et faisons varier leur sens. Le capital va de $6\,858$ à $8\,697$~M\euro{}
aux $1\,024$ sommets admissibles, la largeur variant de $189$~M\euro{} (écart-type) d'une graine
à l'autre ; ces extrema ne sont pas des bornes globales démontrées.
Documenter le sens de six dépendances réduirait cette plage de $32{,}7\,\%$ en moyenne
($22$ à $44\,\%$ selon la graine). Ce scénario précise l'intérêt d'un registre d'incidents,
sans garantir le gain d'une collecte réelle.

\paragraph{Mots-clés.} Risque cyber ; DORA ; SCR ; Solvabilité II ; théorie des valeurs
extrêmes ; cascade dirigée ; identification partielle ; bornes de capital ; inférence
bayésienne ; valeur de l'information.

% USAGE DE L'IA GENERATIVE, VERSION INSTITUT A 150 PAGES. Exige par le § 5.1.c des Recommandations
% Jury de l'Institut des actuaires (« a la fin du resume et distinct de ce dernier »). Doit rester
% coherent avec l'annexe ann:ia.
\bigskip
\noindent\textbf{Usage de l'intelligence artificielle générative.} Des outils d'IA générative
ont servi à la rédaction et à la reformulation du manuscrit, à l'écriture des scripts d'analyse
et à la relecture du texte, sous le contrôle de l'auteur. Les analyses, les choix
méthodologiques et l'interprétation des résultats relèvent de l'auteur. Aucun résultat chiffré
n'est produit par l'assistant : tout nombre publié provient de l'exécution d'un script du dépôt.
L'annexe~\ref{ann:ia} détaille cet usage et les contrôles qui l'encadrent.

% =====================================================================
% USAGE DE L'IA GENERATIVE, VERSION INSTITUT A 150 PAGES. Exige par le § 5.1.c des
% Recommandations Jury de l'Institut des actuaires (« a la fin du resume et distinct de ce
% dernier »). Le texte decrit l'usage reel : redaction et reformulation du manuscrit, scripts,
% relecture. Il doit rester coherent avec l'annexe ann:ia.
% =====================================================================

% << ABSTRACT >> EST UN TITRE, ET NON UNE AMORCE EN GRAS. La liste de controle
% de l'ecole l'exige nommement : << le titre de votre resume francais est
% Resume et le titre de votre resume en anglais est Abstract >>. Il portait
% jusqu'au 28 septembre 2026 un simple  \textbf{Abstract.}  en tete de
% paragraphe, qui n'est pas un titre et n'apparaissait pas au sommaire.
%
% ET LA PAGE NEUVE NE COUTE RIEN. La meme liste demande un resume d'UNE page au
% plus dans chaque langue : le francais et l'anglais tenaient sur la meme page,
% donc l'anglais aurait debordé sur la suivante des qu'on lui pose un titre. Les
% deux blocs occupent desormais une page chacun, et le document n'en gagne
% aucune, le paragraphe d'usage de l'IA occupant deja seul la seconde.
% \chapter* ET NON \section* : le titre de l'abstract doit avoir le meme
% rang que celui du resume francais, la liste de controle les citant cote a
% cote. Un \chapter* ouvre sa page lui-meme, d'ou l'absence de \newpage.
\chapter*{Abstract}
\addcontentsline{toc}{chapter}{Abstract}
\markboth{Abstract}{}

\begin{otherlanguage*}{english}
The European DORA regulation requires financial entities to control their IT risks.
We compare capital needs across assumed compliance states, then examine their sensitivity
to missing information on contagion and costs.

The model combines annual incident counts, losses and propagation across five control areas.
Capital need is the $99.5\,\%$ quantile of annual loss, with costs added across affected pillars.
Regulatory penalties and the cyber policies an insurer sells are excluded.

For the financial sample, capital need rises from $6\,049$ to $20\,188$~M\euro{} between the
compliant and non-compliant states, a factor of $3.34$. These amounts depend on scenario parameters
and on the cost rule. For an entity with $20\,000$~M\$ of assets, the range from $131.5$ to
$183.3$~M\euro{} remains an extrapolated order of magnitude.

The sources cannot calibrate propagation. We hold expert-set link strength fixed and vary
direction. Capital ranges from $6\,858$ to $8\,697$~M\euro{} at the $1\,024$ admissible vertices,
the width varying by $189$~M\euro{} (standard deviation) across seeds; these extrema are not
proven global bounds. Documenting six link directions would narrow
this range by $32.7\,\%$ on average, from $22$ to $44\,\%$ depending on the simulation seed.
This scenario indicates the potential value of an incident register without guaranteeing
the benefit of collecting actual observations.
\end{otherlanguage*}

% =====================================================================
% SOMMAIRE AU NIVEAU DES SECTIONS, et non des sous-sections. Mesure du
% 22 septembre 2026 : en listant les sous-sections il occupait SIX pages, sur un
% document que Kelian veut imperativement sous 200. Rien n'est retire du
% memoire, seule la profondeur du sommaire change ; les sous-sections gardent
% leur numero et leur titre dans le corps.
\setcounter{tocdepth}{1}
% Sommaire en interligne simple (26 octobre 2026, comme la bibliographie) : police inchangee.
% Corps \small depuis le 28 octobre 2026 : les notes de synthese passent a trois pages chacune
% (liste de l'ecole, « entre 3 et 7 pages ») ; le sommaire rend une des deux pages.
\begingroup\singlespacing\small
\makeatletter\patchcmd{\l@chapter}{\vskip 1.0em \@plus\p@}{\vskip 0.15em \@plus\p@}{}{}\makeatother
\tableofcontents
\endgroup

% =====================================================================

% SOURCES-SCRIPTS: 28 39 33 43 68 92 66 89 91 100 30 55 65 117 130 134
